\section*{Capítulo \ref{ch:g12:contdist} --- Variáveis aleatórias contínuas}

\begin{solution}{exo:g12:contdist:1}
$\P(X \leq 5) = \frac{5}{15} = \frac13$;
$\P(4 \leq X \leq 10) = \frac{6}{15} = \frac25$;
$\E(X) = \frac{15}{2} = 7.5$ min;
$\sigma(X) = \frac{15}{\sqrt{12}} = \frac{5\sqrt3}{2} \approx 4.3$ min.
\end{solution}

\begin{solution}{exo:g12:contdist:2}
$f \geq 0$ e
$\int_0^2 \frac38 t^2\,\dd t = \frac38\left[\frac{t^3}{3}\right]_0^2
= \frac38 \times \frac83 = 1$: é uma \omterm{def:g12:contdist:density}{densidade}.
\[
\E(X) = \int_0^2 \frac38 t^3\,\dd t = \frac38 \times \frac{16}{4} = \frac32,
\qquad
\P(X \geq 1) = \int_1^2 \frac38 t^2\,\dd t = \frac{8 - 1}{8} = \frac78 .
\]
\end{solution}

\begin{solution}{exo:g12:contdist:3}
\emph{1.} $\E(X) = \frac{1}{0.2} = 5$ anos;
$\P(X > 5) = \eu^{-0.2\times5} = \eu^{-1} \approx 0.37$.

\emph{2.} Pela \omterm{thm:g12:contdist:memoryless}{ausência de memória}
(\cref{thm:g12:contdist:memoryless}),
$\pcond{X > 3}{X > 8} = \P(X > 5) = \eu^{-1} \approx 0.37$: os três anos
de serviço não mudam nada.
\end{solution}

\begin{solution}{exo:g12:contdist:4}
$\P(X > m) = \eu^{-\lambda m} = \frac12$ dá
$m = \frac{\ln 2}{\lambda} \approx \frac{0.69}{\lambda}$, menor que
$\E(X) = \frac1\lambda$. A \omterm{def:g12:contdist:density}{densidade} exponencial tem uma longa cauda à
direita: alguns tempos de vida atipicamente longos puxam a \emph{\omterm{def:g11:stat:mean}{média}}
acima da \emph{\omterm{def:g11:stat:median}{mediana}}, o valor que metade da população ultrapassa. (Esta
é a meia-vida do \cref{ex:g12:exp:halflife}: metade dos átomos sobrevive a
ela, embora o tempo de vida médio seja maior.)
\end{solution}

\begin{solution}{exo:g12:contdist:5}
$168 = \mu - \sigma$ e $182 = \mu + \sigma$: cerca de $68\%$.
$189 = \mu + 2\sigma$: acima disso fica metade dos
$100 - 95.4 = 4.6\%$ restantes, ou seja, cerca de $2.3\%$.
$154 = \mu - 3\sigma$: cerca de $\frac{100 - 99.7}{2} = 0.15\%$.
\end{solution}

\begin{solution}{exo:g12:contdist:6}
$\P(X < 500) \leq 2.5\%$ significa que $500$ tem de ficar pelo menos dois
\omterm{def:g11:stat:variance}{desvios padrão} abaixo da \omterm{def:g11:stat:mean}{média} (a normal deixa $2.3\% \approx 2.5\%$
abaixo de $\mu - 2\sigma$; com o convencional $1.96$, o raciocínio é
idêntico):
\[
\mu - 2\sigma \geq 500 \iff \mu \geq 500 + 8 = 508 \text{ g}.
\]
A máquina tem de ser ajustada para $508$ g em \omterm{def:g11:stat:mean}{média} --- o preço da
garantia são $8$ g de produto ``de graça'' por pacote.
\end{solution}

\begin{solution}{exo:g12:contdist:7}
\emph{1.} Com $p = \frac16$, $n = 1200$:
$\sqrt{\frac{p(1-p)}{n}} = \sqrt{\frac{5/36}{1200}} \approx 0.0108$, de
modo que o \omterm{def:g10:numbers:interval}{intervalo} de flutuação a $95\%$ é
\[
\frac16 \pm 1.96 \times 0.0108 \approx \intcc{0.146}{0.188}.
\]

\emph{2.} A \omterm{def:g10:stats:series}{frequência} observada $\frac{260}{1200} \approx 0.217$ fica bem
fora: a hipótese de honestidade é rejeitada ao nível de $5\%$. O desvio é
de cerca de $4.6$ \omterm{def:g11:stat:variance}{desvios padrão} --- para uma \omterm{def:g10:numbers:approx}{aproximação} normal, uma
\omterm{def:g10:proba:distribution}{probabilidade} da ordem de $10^{-6}$, muito mais conclusiva que a cota
$\leq 4.6\%$ de Bienaymé--Chebyshev (\cref{exo:g12:sums:7}).
\end{solution}

\begin{solution}{exo:g12:contdist:8}
\emph{1.} A margem $\frac{1}{\sqrt n} \leq 0.02$ exige $n \geq 2500$
pessoas.

\emph{2.} Para distinguir votações separadas por $1$ ponto, a margem tem
de ficar bem abaixo de $0.5$ ponto, o que exige
$n \geq \left(\frac{1}{0.005}\right)^2 = 40\,000$ pela fórmula
simplificada --- já impraticável para a maioria das pesquisas. Pior: a
margem estatística só leva em conta o erro de \emph{amostragem}; os
vieses sistemáticos (\omterm{def:g10:proba:sample}{amostras} não representativas, não resposta, viradas
de última hora) não encolhem quando $n$ cresce e costumam superar um
ponto. Uma disputa de $1$ ponto é genuinamente indecidível.
\end{solution}

\begin{solution}{exo:g12:contdist:9}
\emph{1.} Como $\alpha > 0$, $c \leq \alpha X + \beta \leq d
\iff \frac{c - \beta}{\alpha} \leq X \leq \frac{d - \beta}{\alpha}$, de
modo que
\[
\P(c \leq Y \leq d)
= \int_{(c-\beta)/\alpha}^{(d-\beta)/\alpha} f(t)\,\dd t .
\]
A substituição $y = \alpha t + \beta$ (isto é, ler a área na variável
$y$, com $t = \frac{y - \beta}{\alpha}$, $\dd t = \frac{\dd y}{\alpha}$)
transforma isso em
$\int_c^d \frac{1}{\alpha} f\left(\frac{y-\beta}{\alpha}\right)\dd y$: a
\omterm{def:g11:func:function}{função} $g(y) = \frac1\alpha f\left(\frac{y-\beta}{\alpha}\right)$ é uma
\omterm{def:g12:contdist:density}{densidade} de $Y$.

\emph{2.} Com $f = \varphi$, $\alpha = \sigma$, $\beta = \mu$:
\[
g(y) = \frac{1}{\sigma}\,\varphi\!\left(\frac{y - \mu}{\sigma}\right)
= \frac{1}{\sigma\sqrt{2\pi}}\,
\exp\left(-\frac{(y-\mu)^2}{2\sigma^2}\right),
\]
que é, portanto, a \omterm{def:g12:contdist:density}{densidade} de
$\mathcal N(\mu, \sigma^2) = \mu + \sigma\,\mathcal N(0,1)$.
\end{solution}

\begin{solution}{pb:g12:contdist:1}
\textbf{1.} $\P(\text{espera} \leq 10) = \frac{10}{30} =
\frac13$; $\E = 15$ min;
$\sigma = \frac{30}{\sqrt{12}} \approx 8.7$ min.

\textbf{2.} $\int_0^1 2x\,\dd x = 1$: é uma \omterm{def:g12:contdist:density}{densidade}.
$\E(X) = \int_0^1 2x^2\,\dd x = \frac23$;
$\P\left(X > \frac12\right) = \int_{1/2}^1 2x\,\dd x =
1 - \frac14 = \frac34$.

\textbf{3.} $\P(T > 15) = \eu^{-15/10} \approx 0.22$.
\omterm{def:g11:stat:median}{Mediana}: $\eu^{-m/10} = \frac12$: $m = 10\ln 2 \approx 6.9$
min --- menos que a \omterm{def:g11:stat:mean}{média} $10$, porque a longa cauda à direita
da exponencial (esperas enormes e raras) arrasta a \omterm{def:g11:stat:mean}{média} para
cima, deixando a \omterm{def:g11:stat:median}{mediana} na espera ``típica''.

\textbf{4.} Pela \omterm{thm:g12:contdist:memoryless}{ausência de memória},
$\P(T > 35 \mid T > 20) = \P(T > 15) \approx 0.22$: os vinte
minutos já cumpridos não compram nada --- o fluxo não
envelhece.

\textbf{5.} (a) exponencial (o decaimento radioativo é o
modelo para o qual a \omterm{thm:g12:contdist:memoryless}{ausência de memória} foi inventada); (b)
uniforme em $\intcc{0}{8}$ (um horário com defasagem
aleatória); (c) exponencial; (d) nenhum dos dois: uma lâmpada
gasta tem \emph{mais} \omterm{def:g10:proba:distribution}{probabilidade} de queimar na próxima hora
que uma nova, de modo que
$\P(T > s + t \mid T > s) < \P(T > t)$: o envelhecimento
contradiz o teorema da \omterm{thm:g12:contdist:memoryless}{ausência de memória}.

\textbf{6.} $68\,\%$, $95\,\%$, $99.7\,\%$: os três
\omterm{def:g10:numbers:interval}{intervalos} de referência.

\textbf{7.} $z = \frac{190 - 172}{8} = 2.25$: cauda superior
$\approx 1.2\,\%$.

\textbf{8.} $z = 2$: cerca de $2.3\,\%$ --- mais ou menos uma
pessoa em cada $44$.

\textbf{9.} A especificação fica a $\pm 2\sigma$: cerca de
$95.4\,\%$ passam, $4.6\,\%$ são rejeitados --- ruinoso em
escala. A $\pm 6\sigma$, a taxa de falha cai para cerca de
duas partes por \emph{bilhão}: o seis sigma compra o direito
de produzir em massa sem inspecionar em massa.

\textbf{10.} $\sigma = 5$; com correção de \omterm{def:g12:limcont:continuity}{continuidade},
$\P \approx \P\left(\abs Z \leq \frac{5.5}{5}\right) =
\P(\abs Z \leq 1.1) \approx 0.73$. O teorema alisa a tábua de
Galton do \cref{pb:g11:binom:1}: a escadinha de compartimentos
vira o sino \omterm{def:g12:limcont:continuity}{contínuo}.

\textbf{11.} $n \approx \frac{1.96^2 \times 0.25}{0.005^2}
\approx 38\,400$ pessoas --- e, nesse tamanho, os erros que
\emph{não} são de amostragem (cadastro, recusas, mentiras)
engolem a margem: uma disputa de um ponto está além do
alcance de uma pesquisa honesta, e é por isso que os
institutos sérios então dizem ``empate técnico''.

\textbf{12.} \omterm{thm:g12:contdist:memoryless}{Ausência de memória}: a partir do instante em que
você chega, a espera restante é exponencial de \omterm{def:g11:stat:mean}{média} $10$ ---
os $10$ minutos completos, não $5$. O ``intervalo médio $10$''
da tabela engana: sua espera tem a mesma \omterm{def:g12:randvar:rv}{distribuição} de um
intervalo inteiro.

\textbf{13.} As chegadas aleatórias caem proporcionalmente ao
comprimento do intervalo:
$\P(\text{intervalo longo}) = \frac{15}{20} = \frac34$. Espera
esperada: $\frac34 \times \frac{15}{2} + \frac14 \times
\frac52 = 5.625 + 0.625 = 6.25$ min --- mais que os ingênuos
$5$, embora o intervalo médio seja $10$. Culpada: a
\emph{amostragem enviesada pelo comprimento} --- intervalos
longos apanham mais passageiros.

\textbf{14.} As esperas para trás e para a frente são cada uma
exponencial de \omterm{def:g11:stat:mean}{média} $10$ (a \omterm{thm:g12:contdist:memoryless}{ausência de memória} vale nos dois
sentidos a partir de sua chegada): o intervalo que contém você
tem \omterm{def:g11:stat:mean}{média} de $20$ minutos --- o dobro do intervalo típico.
Paradoxo da inspeção em uma frase: \emph{o intervalo que por
acaso você inspeciona não é um intervalo típico, porque você
tinha mais chance de cair em um intervalo grande.}

\textbf{15.} Turma \omterm{def:g11:stat:mean}{média}: $\frac{9 \times 10 + 110}{10} = 20$
alunos. \omterm{def:g11:stat:mean}{Média} vivida pelos alunos:
$\frac{90 \times 10 + 110 \times 110}{200} = \frac{13\,000}
{200} = 65$ alunos. O folheto imprime $20$; três em cada cinco
alunos sentam na turma de $110$ e vivem os $65$.

\textbf{16.} As pessoas populares aparecem em muitas listas de
amigos, de modo que amostrar ``um amigo'' é enviesado em favor
dos sociáveis --- seus amigos são sorteados nos intervalos
longos da tabela de horários social.

\textbf{17.} Para $t \geq 0$:
$\P(T > t) = \P(-10\ln U > t) = \P\left(U <
\eu^{-t/10}\right) = \eu^{-t/10}$: exatamente a \omterm{def:g11:func:function}{função} de
sobrevivência da exponencial --- um \omterm{def:g12:exp:ln}{logaritmo} converte o ruído
uniforme do computador em qualquer tempo de espera.

\textbf{18.} Cada uniforme tem \omterm{def:g11:stat:mean}{média} $\frac12$ e \omterm{def:g12:randvar:exp}{variância}
$\frac{1}{12}$: a soma de doze tem \omterm{def:g11:stat:mean}{média} $6$ e \omterm{def:g12:randvar:exp}{variância} $1$
--- de modo que a receita devolve \omterm{def:g11:stat:mean}{média} $0$ e \omterm{def:g12:randvar:exp}{variância} $1$.
Doze empurrõezinhos \omterm{def:g12:sums:indep}{independentes} somados: o mecanismo de
Galton e, com a bênção de Moivre--Laplace, o histograma já
tem forma de sino a olho nu.

\textbf{19.} Sobre o modelo: os mercados não são somas de
muitas causas pequenas \emph{\omterm{def:g12:sums:indep}{independentes}} --- os pânicos
correlacionam tudo (a lição de 2008 do problema anterior) e
produzem caudas pesadas que nenhuma \omterm{def:g12:contdist:density}{densidade} normal possui.
Quando os dados sussurram $10^{-137}$, o catecismo do
estatístico responde: \emph{quem é rejeitado é o modelo, não o
dia}.

\textbf{20.} Uniforme: todas as posições iguais, o honesto ``só
conheço os limites''. Exponencial: o risco que nunca
envelhece --- decaimentos, chegadas, cliques. Normal: o sino
democrático de muitas causas pequenas e \omterm{def:g12:sums:indep}{independentes} ---
alturas, erros, \omterm{def:g11:stat:mean}{médias}. A ferroada: o que você \omterm{def:g10:proba:sample}{amostra} é
enviesado pela maneira como você esbarrou nele --- ônibus,
turmas, amigos. E o arco: a criança que contava caixinhas de
suco no ano 6, dez volumes de problemas depois, mediu a
multidão com uma curva; o épsilon, a \omterm{def:g12:integ:area}{integral} e o teorema
central do limite esperam nos volumes de graduação para
explicar \emph{por que} o sino dobra por tudo.
\end{solution}
